% Module 1 section file. Included by ../module1_stellarator_equilibrium_geometry.tex.
\section{Worked conceptual examples}
\label{sec:worked_examples}

\subsection{Example 1: pressure as a flux function}
Suppose a nested magnetic surface is labeled by $s$. Start from force balance,
\begin{equation}
\nabla p=\mathbf J\times\mathbf B.
\end{equation}
Dot with $\mathbf B$:
\begin{equation}
\mathbf B\cdot\nabla p=0.
\end{equation}
Parameterize a field line by $\tau$ using $d\mathbf x/d\tau=\mathbf B$. The pressure derivative along the line is
\begin{equation}
\frac{dp}{d\tau}
=\frac{d\mathbf x}{d\tau}\cdot\nabla p
=\mathbf B\cdot\nabla p
=0.
\end{equation}
Therefore $p$ is constant along that field line. If the line densely covers the surface and $p$ is continuous, all points on the surface have the same pressure, so $p=p(s)$.

\subsection{Example 2: field-line closure}
Let $\iota=2/5$ and use the straight-line solution
\begin{equation}
\theta(\zeta)=\theta_0+\frac{2}{5}\zeta.
\end{equation}
After five toroidal turns, $\Delta\zeta=10\pi$. The poloidal advance is
\begin{equation}
\Delta\theta=\frac{2}{5}(10\pi)=4\pi,
\end{equation}
which is two poloidal turns. The angular coordinates return to their starting values modulo $2\pi$, so the ideal field line closes.

\subsection{Example 3: a rational Alfven resonance surface}
Take a mode with $(m,n)=(2,1)$ and the transform profile
\begin{equation}
\iota(s)=0.35+0.30s.
\end{equation}
The simple parallel wave number vanishes when
\begin{equation}
2\iota(s)-1=0.
\end{equation}
Thus
\begin{equation}
2(0.35+0.30s)-1=0
\quad\Rightarrow\quad
s=0.5.
\end{equation}
The surface $s=0.5$ is the rational surface $\iota=1/2$ for this harmonic. In a more complete shear-Alfven model, coupling, compressibility, and finite-frequency effects determine the actual continuum behavior near this location.

\subsection{Example 4: beta estimate}
Let $p=2.0\times10^5\ \mathrm{Pa}$ and $B=2.5\ \mathrm T$. The magnetic pressure is
\begin{equation}
\frac{B^2}{2\mu_0}
=\frac{(2.5)^2}{2(4\pi\times10^{-7})}
\approx2.49\times10^6\ \mathrm{Pa}.
\end{equation}
The local beta is
\begin{equation}
\beta\approx\frac{2.0\times10^5}{2.49\times10^6}
\approx0.080,
\end{equation}
or approximately $8.0\%$. This does not mean magnetic forces exceed pressure forces everywhere by exactly a factor of $12.5$; equilibrium depends on gradients and curvature, not only local energy-density ratios.

\subsection{Example 5: field-period Fourier variation}
Suppose $N_{\mathrm{fp}}=4$ and a surface contains the term
\begin{equation}
R_{1,1}(s)\cos(\theta-4\zeta).
\end{equation}
Increasing $\zeta$ by one field period, $\Delta\zeta=2\pi/4$, changes the phase by
\begin{equation}
-4\Delta\zeta=-2\pi.
\end{equation}
The cosine is unchanged, confirming fourfold toroidal periodicity.

